% ============================================================
%  Глава 1 — Что робот вообще делает
% ============================================================
\chapter{Робот глазами пользователя}
\label{ch:user_view}

\section{Сценарий «D метров вперёд»}

Открываем веб-дашборд (\filepath{compute\_node/frontend/}), идём на страницу
\texttt{/mps}. Видим окружность с маркером, поля «\textit{дистанция $D$}»
и «\textit{целевая скорость $v_\text{target}$}», кнопку «Run on Robot».

\begin{enumerate}
  \item Двигаем маркер на окружности --- выбираем \textbf{целевой курс}
        $\varphi$ (на сколько градусов отвернуть от «прямо вперёд»).
  \item Вводим $D$ и $v_\text{target}$.
  \item Жмём «Run on Robot».
\end{enumerate}

Робот делает три вещи:

\begin{enumerate}
  \item \textbf{Поворачивается на месте} на угол $\varphi$. Едет на нуле,
        крутится только разворот.
  \item \textbf{Проезжает $D$ метров} с целевой скоростью $v_\text{target}$,
        удерживая курс $\varphi$.
  \item \textbf{Останавливается} и публикует результат: сколько проехал,
        какой был перерегулирующий выброс, как долго «успокаивался».
\end{enumerate}

Всё. Один сценарий, конечный. Никакой карты, никакого глобального плана.

\begin{ideabox}
Книжка построена вокруг \emph{этого} сценария. Каждый следующий вопрос
--- «а как робот\dots» --- разбираем именно для этого сценария. Простота
сценария --- наше учебное преимущество: можно вглядеться в каждую формулу
без отвлечения на другие задачи.
\end{ideabox}

\section{Что нужно роботу, чтобы это сделать}

Распакуем то, что мы только что увидели. У робота должна быть способность
ответить на пять вопросов.

\begin{enumerate}
  \item[\textbf{(Q1)}] \textbf{Где я?}
        Чтобы понять, что проехал $D$ метров, надо знать, на сколько
        я уже сдвинулся. Это \emph{состояние позиции}.
        Об этом --- глава~\ref{ch:state}.

  \item[\textbf{(Q2)}] \textbf{Куда я смотрю?}
        Чтобы повернуться к $\varphi$ и не уехать боком, надо знать
        текущий курс $\theta$. Это \emph{состояние курса}.
        Тоже глава~\ref{ch:state}.

  \item[\textbf{(Q3)}] \textbf{Как мне сдвинуться/повернуться, если я скажу
        моторам «езжай»?}
        Если послать команду $u_v = 0.2$~м/с, робот не сразу станет ехать
        $0.2$~м/с --- моторы инерционны. Это \emph{динамика}.
        Об этом --- главы~\ref{ch:kinematics}, \ref{ch:dynamics}.

  \item[\textbf{(Q4)}] \textbf{Какую команду подать, чтобы достичь цели?}
        Не любой набор команд приведёт к цели хорошо. Нужен \emph{регулятор}.
        Об этом --- главы~\ref{ch:controllers}, \ref{ch:lqr}, \ref{ch:mpc}.

  \item[\textbf{(Q5)}] \textbf{Когда остановиться и как сказать «готово»?}
        Это \emph{логика сценария}, конечный автомат и подсчёт метрик.
        Об этом --- главы~\ref{ch:scenario}, \ref{ch:turn_phase}.
\end{enumerate}

\section{Три слоя проекта}

Робот --- не один компьютер, а три. Они общаются через сеть.

\begin{center}
\begin{tikzpicture}[
  node distance=8mm,
  layer/.style={
    draw, rounded corners=3pt, thick,
    align=center, font=\small,
    minimum width=4.6cm, minimum height=1.7cm,
  },
  flow/.style={-{Stealth[length=2.5mm]}, thick}
]
  \node[layer, fill=blue!8]   (ui)  {\textbf{Browser / React}\\
                                     UI \texttt{/mps}, MatrixEditor,\\
                                     Target picker};
  \node[layer, fill=green!8, right=of ui]  (cn) {\textbf{Compute Node (ноут)}\\
                                     FastAPI, MQTT-broker,\\
                                     idealised simulator};
  \node[layer, fill=orange!8, right=of cn] (pi) {\textbf{Raspberry Pi 4}\\
                                     \texttt{mps\_node.py}, FSM,\\
                                     MPC, motor driver};
  \draw[flow] (ui) -- node[above, font=\scriptsize] {REST + WS} (cn);
  \draw[flow] (cn) -- node[above, font=\scriptsize] {MQTT} (pi);
\end{tikzpicture}
\end{center}

Что где живёт:

\begin{itemize}
  \item \textbf{Browser} (React) --- то, что видит человек: пикер курса,
        поля распределения, графики прогона. Жесть математики тут
        минимум, но есть --- глава~\ref{ch:picker} разберёт.
  \item \textbf{Compute Node} (FastAPI на ноутбуке) --- посредник: принимает
        команды от UI, шлёт их роботу по MQTT, сохраняет историю прогонов.
        Заодно умеет \emph{симулировать} прогон без физического робота
        --- это идеальный $\vect{x}_{k+1} = \mat{A}_d\vect{x}_k + \mat{B}_d\vect{u}_k$
        без шумов. Полезно для UI и юнит-тестов.
  \item \textbf{Pi} (Python, MQTT) --- собственно мозг робота. Подписан
        на топик \texttt{mps/scenario/run}; запускает тик-цикл $50$~Гц,
        внутри которого \emph{вся боевая математика} --- считывание
        одометрии, шаг MPC, выдача команды моторам.
\end{itemize}

\begin{codebox}[title={Где что лежит}]
\begin{tabular}{ll}
\filepath{compute\_node/frontend/src/pages/MpsPage.tsx} & UI страница\\
\filepath{compute\_node/frontend/src/lib/targetAngle.ts} & математика пикера\\
\filepath{compute\_node/dashboard/routers/mps.py}        & REST + WS\\
\filepath{compute\_node/mps\_runner.py}                  & идеальный симулятор\\
\filepath{pi\_nodes/nodes/mps\_node.py}                  & оркестратор на Pi\\
\filepath{pi\_nodes/control/state\_space\_model.py}      & plant ($\mat{A}_d, \mat{B}_d$)\\
\filepath{pi\_nodes/control/mpc\_controller.py}          & MPC ($\mat{Q}, \mat{R}$, $\mat{K}$, QP)\\
\filepath{config.yaml}                                   & секция \texttt{mps:} с матрицами\\
\filepath{matlab/main.m}                                 & оффлайн-синтез матриц\\
\end{tabular}
\end{codebox}

\section{Граничные правила}

Перед тем как нырять в математику, зафиксируем правила игры, которые
не выводятся, а просто заданы дизайном модуля. Если ты увидишь в коде
проверку, которая кажется лишней --- скорее всего, она охраняет одно
из этих правил.

\begin{notebox}
\begin{enumerate}
  \item MPC активен \textbf{только в FSM-состоянии} \texttt{DRIVE\_FORWARD\_MPS}.
        В любом другом (IDLE, SEARCHING, GRABBING, ...) команды моторам
        идут от \emph{другого} контроллера (старый PID), а MPC ничего не
        делает.
  \item Сценарий \textbf{финитный}: достижение цели, тайм-аут, abort или
        ошибка --- любое из четырёх событий завершает прогон, и FSM
        возвращается в IDLE.
  \item Между прогонами можно \textbf{пересобрать} матрицы (MQTT-команда
        \texttt{mps/matrices/set}) --- это и есть «учебный пульт».
        Во время прогона матрицы не меняются.
  \item В состоянии \texttt{DRIVE\_FORWARD\_MPS} \textbf{cmd\_vel} (команда
        моторам) приходит \emph{только} от \texttt{mps\_node}. Голосовые
        команды, детекции мяча и джойстик игнорируются.
  \item Симулятор \textbf{идеальный}: $\vect{x}_{k+1} = \mat{A}_d\vect{x}_k +
        \mat{B}_d\vect{u}_k$ без шумов, без \emph{motor lag}, без проскальзывания.
        Робот --- неидеальный, разница между симом и реальным прогоном
        --- это и есть «ошибка моделирования».
\end{enumerate}
\end{notebox}

\section{Что дальше}

Дальше мы аккуратно ответим на пять вопросов (Q1)--(Q5) по очереди. Сначала
$\to$ кинематика (глава~\ref{ch:kinematics}): что значит «робот едет вперёд»
формально. Потом --- откуда состояние берётся (глава~\ref{ch:state}). Затем
динамика, линеаризация, регуляторы, и наконец --- сам сценарий.

Идём.
